\section{Introduction}
\label{introduction}

% Goal: motivate WHY automated parameter optimization is needed.
% Do not explain PureCLIP internals or the objective formula here.

% ── 1. The biological question ─────────────────────────────────────────
%   • RBPs control post-transcriptional gene regulation. Mapping their
%     RNA binding sites at scale and with high precision is a central
%     goal of modern RNA biology.
%   • eCLIP produces genome-wide binding maps, but the raw output is
%     aligned reads — a peak caller is required to convert these into
%     a discrete set of binding sites.

% ── 2. The parameter sensitivity problem ──────────────────────────────
%   • PureCLIP's output is sensitive to a handful of numeric parameters.
%     Small changes in bandwidth or merge distance can shift the number
%     of called sites by orders of magnitude and substantially alter
%     which genomic regions are recovered.
%   • Crucially, no single parameter set is universally optimal: the
%     best values depend on the RBP, cell line, and antibody efficiency
%     of each individual dataset.

% ── 3. Why manual tuning does not scale ───────────────────────────────
%   • A single PureCLIP run on full-genome data takes hours; evaluating
%     even a modest grid of parameter combinations is therefore
%     impractical without automation.
%   • Manual tuning is also subjective: different researchers apply
%     different heuristics, making results hard to reproduce or compare
%     across datasets and publications.
%   • The ENCODE consortium provides dozens of eCLIP datasets; systematic
%     per-dataset tuning requires an automated, reproducible framework.

% ── 4. This project and research question ─────────────────────────────
%   • We built an automated optimization framework that proposes parameter
%     sets, runs the full pipeline, scores the result against biological
%     quality signals, and iterates — replacing manual trial-and-error.
%   • Two optimizer strategies are compared on identical evaluation
%     infrastructure: an LLM agent (biological reasoning + priors) and
%     Optuna/TPE (systematic Bayesian search).
%   • Research question: for a small bioinformatics parameter space,
%     does LLM reasoning add value over principled numerical optimization?
